\begin{enumerate}

\item Cada larguero está a $y_i = 100$ mm del eje horizontal por el centro de gravedad.

\item La inercia es la suma de $B y_i^2$ de los cuatro largueros.

\item La tensión es $\sigma_i = M_z y_i/I_z$, de compresión en los superiores y de tracción en los inferiores.

\end{enumerate}

\noindent\textbf{Operaciones y resultados intermedios.}

\begin{enumerate}

\item Inercia de la sección idealizada.

\[\fitmath{y_{1}=4\cdot 400\cdot {100}^{2}\approx 1{,}6\times10^{7}\,\mathrm{mm^{4}}}\]

\item Tensión en los largueros extremos.

\[\fitmath{y_{2}=\frac{20000000\cdot 100}{4\cdot 400\cdot {100}^{2}}\approx 125\,\mathrm{MPa}}\]

\end{enumerate}
